% !TEX root = ../main.tex
\section{Conclusion}
\label{sec:conclusion}

In this paper, we studied the iteration complexity of single-loop linearized
augmented Lagrangian methods with a fixed penalty. For the family
\(\calC_\theta\), we established a lower bound of \(\Omega(\epsilon^{-4})\)
that holds for any choice of constant center stepsizes whenever
\(\epsilon(\rho+s)\le c_1\), and we showed that the choice \(\beta=1\)
attains this order. With \(\rho\) fixed and \(s=c_s/\epsilon\) for a
sufficiently large constant \(c_s\), the choice \(\beta=0\) and a suitable
primal center stepsize \(\alpha=\Theta(\epsilon)\) attains
\(\mathcal O(\epsilon^{-3})\).

These results delineate the roles of the penalty, dual regularization, and
center updates within \(\calC_\theta\). The combined penalty \(\rho+s\)
connects penalty selection in ALMs with dual regularization in the reduced
minimax formulation. Adjusting the constant center stepsizes cannot
substitute for the growth of \(\rho+s\) required to attain
\(\mathcal O(\epsilon^{-3})\). This growth can be realized by reducing the
dual regularization coefficient \(1/s\) as the target accuracy decreases,
while keeping the penalty \(\rho\) fixed.

An interesting direction for future research is to extend this analysis to
smoothed gradient descent--ascent methods for general nonconvex--concave
minimax problems.
